\cotefile{29}
\scanpage{33}
% archive sheet number: 32

\uline{Rem\unsure{\deleted{r}o|\deleted{o}r}\unsure{\deleted{o}r|\deleted{r}o}ds} concernant l'exemple 3 du par.1. - La présentation de cet exemple n'est ni élégante ni naturelle, la donnée de $\underline{s}'$ et de \quad étant un élément de structure artificiel. Il vaut mieux présenter les choses ainsi, en termes de conditions sur $\underline{R}$, savoir les suivantes:

a) Pour tout objet $M$ d'un $\underline{C}_{S'}$, le faisceau \inserted{étale} associé au préfaisceau sur $\underline{\operatorname{Et}}_{S'}$ des automorphismes de $M$ est \deleted{d\unsure{u|e}} représentable par un schéma en groupes $\underline{G}(M)$ fini et étale sur $S$ (évidemment déterminé à isomorphisme unique près.

Il est évident que $G(M)$ opère sur le schéma relatif $\unsure{\underline{s}(M)=Z|\underline{s}(M)\equiv Z}$.

b) Pour tout objet $S'$ de $\underline{\operatorname{Et}}_S$, tout morphisme \inserted{\deleted{$M'\quad M$}} dans $\underline{C}_{S'}$ est un épimorphisme, et tout endomorphisme d'un obje\unsure{\deleted{y}t|\deleted{t}y} de $\underline{C}_{S'}$ est un automorphisme. Pour tout morphisme $\inserted{u:}M'\quad M$ dans $\underline{C}_{S'}$, et tout automorphisme $g'$ de $M'$, il existe un automorphisme $g$ de $M$ tel que $gu=ug'$.

Cet automorphisme est évidemment unique, de plus la condition précédente implique que à $\inserted{u:}M'\quad M$ est associé un unique homomorphisme d schémas en groupes $G(M')\quad G(M)$ tel que $u$ soit compatible avec les opér\unsure{a}tions de $G(M')$ et $G(M)$ sur $M'$ et $M$.

c) Avec les notations précédentes, \deleted{si} l'homomorphisme $G(M')\quad G(M)$ est surjectif, et si $H$ désigne son noyau, l\unsure{\deleted{'}e|\deleted{e}'}\deleted{homo}morphisme $\underline{s}(M')\quad\underline{s}(M)$ induit un \deleted{morphis\unsure{m}e} isomorphisme $\underline{s}(M')/H\quad\underline{s}(M)$, \deleted{i\unsure{.}e\unsure{.} fait de $\underline{s}(M')$ un \unsure{fibéé|fib\'{r}é} principal sur $\underline{s}(M)$, de groupe} (NB on suppose que les $\underline{s}(M)$ sont \deleted{\unsure{relativement} des \unsure{revêtem}} relativement affines, de sorte que le passage au quotient a un sens).

d) Deux homomorphismes de $M'$\deleted{d} dans $M$ sont,localement pour la top\unsure{\deleted{I}o|\deleted{t}o|\deleted{l}o|\deleted{o}I|\deleted{o}t|\deleted{o}l}logie étale, déduits l'un de l'autre par composition avec un automorphi\unsure{s} de $M$.

